\section{\sys{} 方法}
\label{sec:method}
% 跟着系统图走。每一步依次写目的、定义、算法、保证、复杂度。
\sys{} 由两步组成。第 1 步<……>。第 2 步<……>，其输出反馈给下一轮的第 1 步。

\subsection{第 1 步：<名称>}
\label{sec:step1}
<一句话说明目的>。<定义>。<在贯穿例子上的演示>。

\subsection{第 2 步：<名称>}
\label{sec:step2}
<一句话说明目的>。算法~\ref{alg:main}列出了具体过程。

\begin{algorithm}[t]
\caption{<过程名称>}
\label{alg:main}
\begin{algorithmic}[1]
\REQUIRE 候选集 $\cand$，预算 $\budget$
\ENSURE 已标注集合 $L$
\STATE $L \gets \emptyset$
\WHILE{$|L| < \budget$}
  \STATE <按第 2 步的规则选择下一个候选>
\ENDWHILE
\RETURN $L$
\end{algorithmic}
\end{algorithm}

\begin{theorem}[保证]
\label{thm:guarantee}
在假设 <k> 下，算法~\ref{alg:main}的策略不低于<它所包含的每个基线>。
\end{theorem}
% 实现不满足定理的全部前提时，把定理写成条件保证。
